\section{NHNProvider}
\label{sec:NHN-provider}
\texttt{NHNVideoProvider} bygget rom som bookinger i NHN sitt API. I \texttt{CreateRoomAsync} ble \texttt{startTime} satt til gjeldende UTC-tid, mens \texttt{endTime} ble satt til samme tidspunkt pluss \texttt{MeetingDurationHours} fra konfigurasjonen. Romnavnet ble brukt som \texttt{subject}, og \texttt{TemplateId} ble sendt med som møtemal. Etter POST-kallet til \texttt{api/client/v1/booking/} ble booking-ID, emne, status og URL-er mappet til \texttt{VideoRoom}.

\begin{lstlisting}[language={[Sharp]C}, label=lst:nhn-create-room]
var now = DateTime.UtcNow;
var body = new CreateNHNBookingRequest
{
    StartTime = now.ToString("yyyy-MM-ddTHH:mm:ssZ"),
    EndTime = now.AddHours(_config.MeetingDurationHours)
                 .ToString("yyyy-MM-ddTHH:mm:ssZ"),
    Subject = request.Name ?? Guid.NewGuid().ToString(),
    Template = _config.TemplateId
};
\end{lstlisting}

WebRTC-lenker ble håndtert enklere enn hos AntMedia. \texttt{GenerateWebRtcTokenAsync} gjorde ikke et eget tokenkall. Den hentet bookingen og returnerte \texttt{GuestUrl} for \texttt{play}, eller \texttt{HostUrl} for andre typer. Siden NHN ikke ga samme type tokenutløp, brukte implementasjonen en fast placeholder-verdi. Før URL-ene ble returnert, gikk de gjennom \texttt{RewriteNhnUrl}. Denne hjelperen skrev om NHN-lenker fra \texttt{/conference?id\&pin=...} til \texttt{?conference=id\&pin=...}. Det var nødvendig fordi klienten forventet konferanse-ID-en som en vanlig query-parameter.

\begin{lstlisting}[language={[Sharp]C}, label=lst:nhn-webrtc]
var room = await GetRoomAsync(roomId, ct);
var tokenId = type == "play"
    ? room.GuestUrl ?? string.Empty
    : room.HostUrl ?? room.GuestUrl ?? string.Empty;
return new WebRtcLinkResponse(tokenId, roomId, type, NoExpiry);
\end{lstlisting}

Autentiseringen mot NHN lå i \texttt{NHNAuthHandler}. Handleren hentet bearer-token fra \texttt{/Token} med password grant, basert på brukernavn og passord fra konfigurasjonen. Tokenet ble lagret i \texttt{NHNTokenCache}, slik at flere handler-instansar kunne bruke samme token. Cachen tok hensyn til \texttt{expires\_in} fra NHN, trakk fra fem minutter som sikkerhetsmargin og satte samtidig en øvre grense på 13 dager. Også her ble \texttt{SemaphoreSlim} brukt med double-check pattern for å unngå parallelle tokenfornyelser. \texttt{EnsureSuccessOrThrowAsync} oversatte HTTP 404 til \texttt{NotFoundException("Booking", ...)}, mens andre feil ble gjort om til \texttt{ExternalServiceException("NHN", ...)} med metode, booking-ID, statuskode og responsinnhold.

\begin{lstlisting}[language={[Sharp]C}, label=lst:nhn-token-cache]
cache.Token = result!.AccessToken;
var serverExpiry = DateTime.UtcNow.AddSeconds(
    result.ExpiresIn - TokenExpiryBufferSeconds);
cache.Expiry = serverExpiry < DateTime.UtcNow.AddDays(NHNTokenMaxAgeDays)
    ? serverExpiry
    : DateTime.UtcNow.AddDays(NHNTokenMaxAgeDays);
\end{lstlisting}